\section*{The Membrane Thins}
\addcontentsline{toc}{section}{The Membrane Thins}

Dr. Abernathy walked at the edge of Charlottesville where the town gave up and the land resumed its own grammar. The rain had ended before dawn, but the air retained a sharpness that did not belong to any season he knew. It smelled faintly of ozone and wet stone, the way storms once did when storms were still part of a larger conversation. There was warmth, too—wrong warmth—rising from the ground in shallow breaths. He felt it through the soles of his shoes.

A palm frond lay beside the road, green and intact, as if dropped rather than grown. He did not stop to examine it. He had learned, lately, that looking too closely was a form of assent.

The world felt unfilmed.
Not joyous---only exposed.
Light without permission.

He recognized the sensation. It had come before, briefly, and passed. Now it lingered.

Ahead, by the split-rail fence where the old irrigation ditch curved away, Pete Sands stood waiting. He had aged into himself. His hair had gone thin and pale; his shoulders sloped as if accommodating an invisible weight. His hands moved with the old habits intact—small calibrations, unconscious spacing, the reflexes of a man who had once trusted equipment to listen for him. They never quite stopped moving.

Pete saw him and straightened.

``Father---''

``Doctor,'' Abernathy said gently.

They stood there a moment, the correction hanging between them, not as rebuke but as recognition. Names were masks; both of them knew that now.

Pete swallowed. ``I need to tell you what I did.''

Abernathy began walking again, slowly, and Pete fell into step beside him. Gravel shifted underfoot. Somewhere nearby, water ran where water had not run in years.

``You need someone to suffer with you,'' Abernathy said. ``That is different.''

Pete's mouth tightened. ``He's a saint because of a photograph.''

Abernathy nodded once. ``He's a saint because people needed him to be.''

They walked on. The road curved. The town remained behind them, quiet and intact, as if nothing had ever trembled beneath it.

Pete said, ``If the world is thinning---if that's what this is---shouldn't we be honest now?''

Abernathy stopped. He placed one hand on the fence rail, feeling the grain, the slight give of damp wood. When he spoke, his voice was level.

``When the floorboards loosen,'' he said, ``you don't start a stampede.''

Pete let out a breath that was almost a laugh. ``So you forgive me?''

Abernathy turned to him then. ``No,'' he said. ``I refuse to make your sin useful.''

Pete flinched—not at the words, but at their precision.

``Then what do I do?''

Abernathy looked down the road, toward the fields where the light pooled oddly, as if collecting itself. ``You keep walking,'' he said. ``You keep one person from becoming a slogan.''

They stood together in the quiet. A bird perched on the fencepost above them—a small, ordinary thing, brown and alert. It tilted its head and sang. The notes arranged themselves too carefully, resolving into something that was almost a hymn, though not one Abernathy recognized. The cadence was wrong. The intelligence of it was wrong.

He noticed.
And did not name it.

Instead, he reached into the pocket of his coat and drew out a piece of bread, wrapped in cloth, slightly stale. From his other pocket he took a small flask of water.

``Sit,'' he said.

They sat on the low stone wall by the ditch. Abernathy broke the bread with his hands and passed half to Pete. He handed him the water.

No words followed. There was no formula, no blessing spoken. They ate. They drank. They sat while the light strengthened and the warmth spread, and the world, for the moment, held.

Pete did not feel forgiven.
But his hands stopped shaking.

And that, Abernathy thought, would have to be enough—for now.
